\documentclass{article}
\usepackage{../style}
\usepackage{tabularx}

\usepackage[margin = 1.25in]{geometry}

\captionsetup{justification = centering}

\title{Candidacy Presentation Notes}
\author{Zack Dooley}
\date{}

\begin{document}
\maketitle
Syntax - Formal specification of a theory ---- Semantics - Interpretation of a theory in a model
We have a theory of groups - set G, with binary operation $\circ$, satisfying...

This theory can be modeled within $\Z$ w/ addition (externally defined object which satisfies the rules of the theory)


\newpage
\begin{tabularx}{\columnwidth}{X|X}
    Syntax - Formal specification of a theory & Semantics - Interpretation of a theory in a model\\
    \hline
    A Group is a set with binary operation that is unital, associative, has inverses & $\Z$ with addition (this is ONE model, in particular $\Z$ comm. but that is not true of all groups)\\
    (Classical) Prop. Logic: Props. P, Q, etc. $\vee,\wedge,\implies,\neg$, rules about how to use these symbols to create formulas & Truth tables (lets us evaluate formula without constructing a proof)\\
    Dependent type theory (w/ ext. Id) (a class of theories) & LCCC (Clairambault, Dybjer 2011)\\
    Dependent type theory w/ $\Id,1,\Sigma,\Pi_{\ext}$ & $\text{LCCC}_\infty$
\end{tabularx}

Part 1 Syntax:
\begin{itemize}
    \item What is a dependent type theory? A system of rules for deriving valid formula
    \item These rules occur in four judgment forms:
    $$ \Gamma\vdash A\text{ Type},\quad\Gamma\vdash A\equiv B\text{ Type},\quad\Gamma\vdash a:A,\quad\Gamma\vdash a\equiv b:A $$
    (Note distinction between types/terms, terms live in unique types $0:\Z$ not the same as $0:\N$)
    \item Contexts are finite lists of types $\Gamma=(x_1:A_1,\dots,x_n:A_n)$
    \item We call this dependent type theory, because the types/terms can depend on previously occurring variables
    \item Ex: $n:\N\vdash\text{Vec}(n)\text{ Type}$ $n:\N\vdash[]:\text{Vec}(n)$ (Vec(n) is a family of types dependent on $\N$)
    \item (Example of derivation in type theory - variable change rule?)
    \item (there are a set of basic rules which I am ignoring, they're "expected" rules about the basic functionality of the theory)
    \item We add more structures to our theory to reason about more things: 1 (unit type - type with unique term), $\Sigma$ (Sigma/dependent pair type) $(a,b):\Sigma_{x:A}B(x)$ corresponds to $a:A$ and $b:B(a)$, $\Pi$ (Pi/dependent function type) $f:\Pi_{x:A}B(x)$ evaluates to $f(a):B(a)$, Id (identity/path type - not the same as judgmental eq.) if $a,a':A$, then we have a type $\Id_A(a,a')$ (type of paths/equalities between $a$ and $a'$, equalities are not unique, leads to higher structures)
\end{itemize}

Part 2 Semantics:
\begin{itemize}
    \item Simplicial sets?
    \item There are many different models for the "intuitive" notion of an infinity category, the one we're interested in is that of a quasi-category
    \item \begin{definition}
        A quasi-category is a simplicial set $X$, which satisfies the following lifting property,
        \[\begin{tikzcd}
            \Lambda_n^k \arrow[r] \arrow[d] & X\\
            \Delta_n \arrow[ur, dashed]
        \end{tikzcd},\]
        for all $n\in\N$ and $0<k<n$. (Example with $\Lambda^1_2\to\Delta_2$ - how this gives us composition)
    \end{definition}
    \item This is an $(\infty,1)$-category i.e. Higher cells are invertible
    \item uniqueness -> contractibility
    \item Basically any categorical construction can be replicated with higher categories, over categories, terminal object, limits
\end{itemize}

Part 1+2 Conjecture:


\end{document}